% Limitations & Epistemic Scope subsection (≤500w).
% For BioSci's locked slot in the Discussion. Cites Math Lemma S1-exchangeability as a *bound*.
% Covers: exchangeability failure modes, reference-class scope, non-ontic status, circularity mitigations.

\subsection*{Limitations and epistemic scope}
REAL relies on an exchangeability assumption between the known rare populations we hold out for calibration (pancreatic $\epsilon$ cells, pulmonary ionocytes, LAMP3$^+$ mregDC, TPEX, FOLR2$^+$ perivascular macrophages and the other populations enumerated in Table~S\ref{tab:heldout}) and the genuinely unannotated rare subpopulations to which we transfer the estimate. The assumption is formally stated in Lemma~\ref{lem:exchangeability} as a \emph{bound}: the held-out-rare loss rate upper-bounds the expected loss rate for unannotated rare subpopulations of comparable abundance, local geometry, and batch profile. We emphasise four boundary conditions under which the bound is not tight. First, the bound degrades when the unannotated population lies outside the span of the highly-variable genes used to construct the latent space; REAL cannot protect what the representation cannot see. Second, the bound degrades when the unannotated population has a qualitatively different batch-effect signature from the calibration populations---for instance, a disease-specific state observed in only one batch. Third, the bound is conservative for populations whose local density profile is not well-approximated by the level-set stratification we use; transitional cells along a continuum can appear as low-density but are not rare in the sampling sense. Fourth, our use of known rare populations as proxies for unknown ones inherits a selection bias: the populations we can hold out are those that have already been discovered, and the discovered set may not be exchangeable with the undiscovered set.

Our strongest defensible claim is therefore deliberately limited. REAL provides (i) a label-free, invariant-based detectability framework calibrated on held-out known rare subpopulations, (ii) an integration strategy (RareShield) that demonstrably reduces the held-out-rare loss rate at pan-cancer-atlas scale without degrading batch correction for abundant cell types, and (iii) an explicit exchangeability bound under which the held-out loss rate is an upper bound on the loss rate of truly unannotated rare subpopulations. A REAL-flagged motif $w \in \mathcal{W}$ is a claim that reproducible structure has been destroyed, not a claim that a new cell type exists; upgrading to an ontological claim requires the usual experimental burden (sorting, orthogonal markers, spatial validation, replication). The complementary negative: a dataset with low $|\mathcal{W}|$ under a candidate integration is evidence that no cross-batch-reproducible rare structure was erased, but not a guarantee that no such structure existed---by construction, we cannot rule out subpopulations whose reproducibility signal was below our sampling resolution or outside the evaluated feature span. These boundaries are not weaknesses of REAL specifically; they are the price of reasoning about what the evaluation framework cannot directly observe. Naming the price is what distinguishes an epistemically calibrated claim from a framework-captured one.
